
%--------------------------------------------------------
\subsection{Finite-Dimensional Factor Representations}
%--------------------------------------------------------

A general term structure model describes an entire curve of bond prices across maturities. In practice, however, it is often assumed that the term structure is driven by a finite-dimensional state vector
\[
z_t \in \mathbb{R}^{\ell}.
\]

Under this assumption, zero-coupon bond prices admit the representation
\begin{equation}
P(t,T)=P(z_t,\tau), \qquad \tau=T-t.
\label{eq:state_representation}
\end{equation}
Hence, the full term structure at time \(t\) is determined by the current value of the state vector \(z_t\).

The short rate may then be written as a function of the state variable,
\begin{equation}
r(t)=\tilde r(z_t),
\label{eq:short_rate_state}
\end{equation}
so that the instantaneous risk-free rate is also determined by the current state.

Since bond prices are characterized by arbitrage-free pricing, the relevant state dynamics for valuation are the dynamics under the risk-neutral measure. We therefore assume that the state process evolves according to
\begin{equation}
dz_t=\mu(z_t)\,dt+\sigma(z_t)\,dW_t^{\mathbb Q},
\label{eq:state_dynamics}
\end{equation}
where \(\mu(\cdot)\) and \(\sigma(\cdot)\) denote the drift and diffusion functions under \(\mathbb Q\).

This representation provides a tractable way of modeling the term structure through a small number of underlying factors rather than as an unrestricted infinite-dimensional object. It also forms the basis for many classical term structure models, including affine term structure models.




Correlation theory:





The setup above assumes independent Brownian motions. \textcolor{red}{In many applications (måske specificere til denne opgave)}, however, the driving disturbance terms are correlated, and we therefore extend the framework to accommodate this. Given the independent Brownian motions $\bar{W}^1,...,\bar{W}^d$ and a deterministic constant matrix
$$\delta = \begin{bmatrix}
\delta_{11} & \delta_{12} & \dots  & \delta_{1d} \\
\delta_{21} & \delta_{22} & \dots  & \delta_{2d} \\
\vdots      & \vdots      & \ddots & \vdots      \\
\delta_{n1} & \delta_{n2} & \dots  & \delta_{nd}
\end{bmatrix}$$
we define the $n$-dimensional process $W=(W_1,...,W^n)^\top$ by $W=\delta \bar{W}$, or componentwise
$$W_t^i=\sum_{j=1}^d \delta_{ij}\bar{W}_t^j,\;\;\;i=1,\dots,n$$

Assuming that the rows of $\delta$ have unit length, i.e.
$$||\delta_i||=\sqrt{\sum_{j=1}^d \delta_{ij}^2}=1,\;\;\;\;\;i=1,...,n,$$
each component $W^i$ is itself a standard Brownian motion with unit variance. The instantaneous correlation structure between the components is captures by the correlation matrix $\rho$, defined by
$$\rho_{ij}\,dt=\text{Cov}[dW_t^i,dW_t^j],$$
which can be shown to satisfy $\rho=\delta\delta^\top$. The process $W$ is then called a vector of correlated Brownian motions with correlation matrix $\rho$.


\begin{definition}[Correlated Brownian Motions]
The process $W=(W^1,...,W^n)^\top$, constructed as above, is called a vector of correlated Brownian motions with correlation matrix $\rho=\delta\delta^\top$.
\end{definition}

The multiplication table from the independent case stated in Remark \ref{rem:multitable} is then modified, the only change being that
$$dW_t^i\cdot dW_t^j=\rho_{ij}\,dt,$$
replacing the earlier rule $dW_t^i\cdot dW_t^j=0$ for $i\neq j$.


THe old subsection of the generator

\subsection{The Infinitesimal Generator}


Examining the drift term in Itô's formula, we observe that the dynamics $f(t,X_t)$ are governed by two components: the explicit time derivative $\partial_t f$ and a purely spatial part determined by the coefficients $\mu$ and $\sigma$ of the SDE \eqref{eq:SDE}. This spatial part appears repeatedly in the analysis and it is therefore convenient to give it a name.

\begin{definition}[Infinitesimal Generator]
Given the SDE defined in \eqref{eq:SDE}, the infinitesimal generator $\mathcal{P}$ is defined for any $f\in C^2(\mathbb{R}^n)$ by
\begin{equation}
    \mathcal{P}f(x) = \sum_{i=1}^n \mu_i(t,x) \frac{\partial f}{\partial x_i}(x) + \frac{1}{2} \sum_{i,j=1}^n (\sigma\sigma^\top)_{ij}(t,x) \frac{\partial^2 f}{\partial x_i \partial x_j}(x).
\end{equation}
The operator $\mathcal{P}$ differentiates with respect to $x$ only with time entering through the coefficients $\mu(t,x)$ and $\sigma(t,x)$.
\label{def:generator}
\end{definition}

In terms of $\mathcal{P}$, Itô's formula for a time-dependent function $f(t,X_t)$ takes the compact form 
\begin{equation}
    df(t, X_t) = \left(\partial_t f + \mathcal{P}f\right)dt + \nabla_x f \cdot \sigma\, dW_t.
\end{equation}
In particular, recalling Corollary \ref{cor:SImartingale}, $f(t,X_t)$ is a martingale if and only if $\partial_t f+\mathcal{P}f=0$, which is a partial differential equation (PDE) in $(t,x)$. \textcolor{red}{This connection between the probabilistic dynamics of $X$ and the theory of PDEs is the foundation of analysis carried on later on.}



Fredes foreslag til itos formel afsntitet med ændret struktur ift correleret BMs:

\subsection{Itô's Formula}

Having established the stochastic integral, we now ask how to 
differentiate a smooth function $f(t, X_t)$ of a diffusion process. 
In ordinary calculus, the chain rule gives 
$df = \partial_t f\,dt + \nabla_x f\,dX$. However, this fails for 
stochastic processes: because Brownian motion has non-zero quadratic 
variation, $(dW_t)^2 = dt$, the second-order term in the Taylor 
expansion does not vanish and must be retained. Itô's formula is the 
stochastic analogue of the chain rule that accounts for this correction.

We work in the general setting of an $n$-dimensional process 
$X = (X^1, \ldots, X^n)^\top$ driven by $d$ Brownian motions 
$W = (W^1, \ldots, W^d)^\top$ with instantaneous correlation structure
\begin{equation}
    dW^i_t \cdot dW^j_t = \rho_{ij}\,dt,
\end{equation}
where $\rho$ is the correlation matrix. The dynamics of $X$ are given by
\begin{equation}
    dX_t = \mu_t\,dt + \sigma_t\,dW_t,

\end{equation}
where $\mu_t \in \mathbb{R}^n$ is the drift vector and 
$\sigma_t \in \mathbb{R}^{n \times d}$ is the diffusion matrix.

\begin{theorem}[Itô's Formula]
Let $X$ have dynamics given by \eqref{eq:multiSDE} and let 
$f \in C^{1,2}(\mathbb{R}_+ \times \mathbb{R}^n)$. Then $f(t,X_t)$ 
is itself an Itô process with stochastic differential
\begin{equation}
    df(t,X_t) = \left\{
        \frac{\partial f}{\partial t}
        + \sum_{i=1}^n \mu_t^i \frac{\partial f}{\partial x^i}
        + \frac{1}{2}\,\mathrm{tr}\!\left[\sigma_t^\top H_f \sigma_t\right]
    \right\}dt
    + \sum_{i=1}^n \frac{\partial f}{\partial x^i}\,\sigma_t^i\,dW_t,
\end{equation}
where $H_f$ denotes the Hessian with entries 
$[H_f]_{ij} = \partial^2 f / \partial x^i \partial x^j$, and $\sigma_t^i$ 
denotes the $i$-th row of $\sigma_t$. Compared to the ordinary chain 
rule, the additional term $\frac{1}{2}\,\mathrm{tr}[\sigma_t^\top H_f 
\sigma_t]\,dt$ arises precisely from the non-zero quadratic variation 
of $W$, and is the defining correction of Itô calculus. The formula 
follows from the formal multiplication rules
\begin{equation}
\begin{cases}
    (dt)^2 = 0,\\
    dt \cdot dW^i_t = 0,\\
    dW^i_t \cdot dW^j_t = \rho_{ij}\,dt,
\end{cases}
\end{equation}
where the last rule reduces to $dW^i_t \cdot dW^j_t = 0$ for $i \neq j$ 
in the independent case.
\end{theorem}

\begin{remark}[Infinitesimal Generator]
The drift in \eqref{eq:Ito} decomposes naturally into an explicit time 
derivative $\partial_t f$ and a purely spatial part determined by 
$\mu$ and $\sigma$. Defining the latter as the infinitesimal generator
\begin{equation}
    \mathcal{A}f(x) = \sum_{i=1}^n \mu_i\frac{\partial f}{\partial x_i}
    + \frac{1}{2}\sum_{i,j=1}^n (\sigma\sigma^\top)_{ij}
    \frac{\partial^2 f}{\partial x_i \partial x_j},
\end{equation}
Itô's formula takes the compact form 
$df = (\partial_t f + \mathcal{A}f)\,dt + \nabla_x f \cdot \sigma\,dW_t$.
In particular, since the stochastic integral has zero mean by 
Proposition~\ref{prop:SIproperties}, $f(t,X_t)$ is a martingale if and 
only if $\partial_t f + \mathcal{A}f = 0$. This connection between the 
dynamics of $X$ and the theory of PDEs will resurface in 
Section~\ref{sec:pricing}.
\end{remark}



Old plan for Cauchy problem generalisation:


\subsection{Placeholder for Cauchy problem for now}

Now we know what the SDE is, when a solution exists, but the biggest question of them all, is how to find it. This is where Feynman Kac comes into the picture. However in a very specific setup of Cauchy problems that we will now introduce.


Still considering the $n$-dimensional SDE, we are in need of \textcolor{red}{write the motivation for looking at the operator theory here!}

Through the Itô formula, the process above is closely connected to a \textbf{partial differential operator} $\mathcal{P}$, which we define below

\begin{definition}
    Given the SDE in \textcolor{red}{reference to the sde above} the partial differential operator $\mathcal{P}$, referred to as the infinitesimal operator of $X$, is defined, for any function $h$ with $h\in C^2(\mathbb{R}^n)$, by

    \begin{equation}
        \mathcal{P} h(t,x)= \sum_{i=1}^{n} \mu_i(t,x) \frac{\partial h}{\partial x_i} (x) + \frac{1}{2} \sum_{i,j=1}^{n} (\sigma \sigma^\top)_{ij}(t,x) \frac{\partial^2 h}{\partial x_i \partial x_j}(x)
    \end{equation}

    \begin{equation}
        \mathcal{P} = \mu(t,x)\frac{\partial}{\partial x} + \frac{1}{2}\sigma^2(t,x)\frac{\partial^2}{\partial x^2}
    \end{equation}
\end{definition}
\textcolor{red}{why does h depend only on x in some parts?}

we note that in terms of the infinitesimal generator, the Itô formula takes the form
... (page 72)


This is an important result in order to define the problem we intend to selv with the modeling of yield curves, which we will start of formulating in a purely mathematical sense and afterwards.

Now we know what the SDE is, when a solution exists, but the biggest question of them all, is how to find it. This is where Feynman Kac comes into the picture. However in a very specific setup of Cauchy problems that we will now introduce.

We will define and example of the so called \textbf{Cauchy Problem} which is a \textbf{boundary value problem}


\textcolor{red}{so we need define the problem that is what we will essentially solve}

and then how Feynman Kac is way to find one of the "nice" solutions.





From the article FredNilsLuca we have the more general:

We start by defining the parabolic Cauchy Problem from \textcolor{red}{FredNilsLuca}, which follow Friedman \textcolor{red}{(1975)}.

We are interested in solving the following backward parabolic Cauchy problem in $\mathbb{R}^n$

\textcolor{write up there problem but in our framework of pricing perspectiv, but still purelty matehatmical. Then after the ML section, we can introduce the final problem, that we intent to solve with neural networks and hereof write about the benchmarks in the same section.}





% infinisteaml cases:


\textcolor{red}{
\begin{remark}[Infinitesimal Generator \textcolor{red}{(1dim)}]
The drift in \eqref{eq:Ito} decomposes naturally into an explicit time 
derivative $\partial_t f$ and a purely spatial part determined by 
$\mu$ and $\sigma$. Defining the latter as the infinitesimal generator
\begin{equation}
    \mathcal{A}f(x) = \sum_{i=1}^n \mu_i\frac{\partial f}{\partial x_i}
    + \frac{1}{2}\sum_{i,j=1}^n (\sigma\sigma^\top)_{ij}
    \frac{\partial^2 f}{\partial x_i \partial x_j},
\end{equation}
Itô's formula takes the compact form 
$df = (\partial_t f + \mathcal{A}f)\,dt + \nabla_x f \cdot \sigma\,dW_t$.
In particular, since the stochastic integral has zero mean by 
Proposition~\ref{prop:SIproperties}, $f(t,X_t)$ is a martingale if and 
only if $\partial_t f + \mathcal{A}f = 0$. This connection between the dynamics of $X$ and the theory of partial differential equations is made precise by the Feynman-Kac formula of Section~\ref{subsec:fk}, and underlies the pricing framework developed in later chapters.
\end{remark}
}

\textcolor{red}{
\begin{remark}[Infinitesimal Generator \textcolor{red}{(multidim)}]
The drift in \eqref{eq:Ito} decomposes naturally into an explicit 
time derivative $\partial_t f$ and a purely spatial part determined by 
$\mu_t$ and $\sigma_t$. Defining the latter as the infinitesimal generator
\begin{equation}
    \mathcal{A}f(x) = \sum_{i=1}^n \mu_t^i \frac{\partial f}{\partial x^i}
    + \frac{1}{2}\operatorname{tr}\!\left[\sigma_t^\top H_f \sigma_t\right],
\end{equation}
where $H_f$ denotes the Hessian matrix with entries 
$[H_f]_{ij} = \frac{\partial^2 f}{\partial x^i \partial x^j}$, 
Itô's formula takes the compact form
\begin{equation}
    df = \left(\partial_t f + \mathcal{A}f\right)dt 
    + \sum_{i=1}^n \frac{\partial f}{\partial x^i}\sigma_t^i\, dW_t.
\end{equation}
In particular, $f(t, X_t)$ is a martingale if and only if 
$\partial_t f + \mathcal{A}f = 0$, which is a partial differential 
equation in $(t,x)$ whose coefficients are exactly those of the 
SDE \eqref{eq:multiSDE}. This connection between the dynamics of $X$ 
and the theory of partial differential equations is the foundation 
of the subsequent analysis.
\end{remark}
}






% OLD FEYNMAN KAC SECTION






\newpage

\subsection{The Feynman--Kac Formula}
\label{subsec:fk}

Remark~\ref{rem:generator} established one direction of the link between diffusion processes and partial differential equations: if $f(t,X_t)$ is a martingale, then its drift must vanish, which leads to the equation
\[
\partial_t f + \mathcal{A}f = 0.
\]
In financial applications, the converse direction is particularly important. Prices are typically expressed as conditional expectations under a pricing measure, and one then seeks the differential equation satisfied by that expectation. For Markov diffusion processes, the Feynman--Kac formula provides precisely this connection between conditional expectations and parabolic PDEs.

We state a version that allows for a zeroth-order term in the PDE. This is the form relevant for asset-pricing applications, where the quantity of interest is a discounted conditional expectation and the discounting enters through the additional term $-k\,u$.

\begin{theorem}[Feynman--Kac, {\citealp[Prop.~5.6]{Bjork2020}}]
\label{thm:fk}
Let $\varphi : \mathbb{R}^n \to \mathbb{R}$ be continuous and satisfy a polynomial growth condition, let $k : [0,T] \times \mathbb{R}^n \to \mathbb{R}$ be bounded and continuous, and suppose that
\[
u \in C([0,T]\times\mathbb{R}^n)\cap C^{1,2}([0,T)\times\mathbb{R}^n)
\]
solves the terminal-value problem
\begin{equation}
\begin{cases}
    \partial_t u(t,x)+\mathcal{A}u(t,x)-k(t,x)\,u(t,x)=0,
    & (t,x)\in[0,T)\times\mathbb{R}^n, \\[4pt]
    u(T,x)=\varphi(x),
    & x\in\mathbb{R}^n,
\end{cases}
\label{eq:fk_pde}
\end{equation}
where $\mathcal{A}$ is the infinitesimal generator defined in \eqref{eq:generator}. Then $u$ admits the probabilistic representation
\begin{equation}
u(t,x)
=
\mathbb{E}\!\left[
    \exp\!\left(-\int_t^T k(s,X_s)\,ds\right)\varphi(X_T)
    \,\middle|\,
    X_t=x
\right],
\label{eq:fk_formula}
\end{equation}
where $X$ is the diffusion process with dynamics \eqref{eq:multiSDE}.
\end{theorem}

\begin{proof}
See \citet[Prop.~5.6]{Bjork2020}. A version under weaker regularity conditions is given in \citet[Thm.~5.3]{Friedman1975}.
\end{proof}

Two aspects of \eqref{eq:fk_formula} are particularly important. First, the representation depends on the past only through the current state $(t,x)$, reflecting the Markov property from Definition~\ref{def:markov}. Second, if $k\equiv 0$, then \eqref{eq:fk_formula} reduces to
\[
u(t,x)=\mathbb{E}\!\left[\varphi(X_T)\mid X_t=x\right],
\]
while \eqref{eq:fk_pde} reduces to
\[
\partial_t u+\mathcal{A}u=0,
\]
which is exactly the drift condition from Remark~\ref{rem:generator}. The Feynman--Kac formula can therefore be viewed as the natural extension of the martingale relation to the case where the conditional expectation is weighted by an exponential discount factor.

This result will be central in the pricing theory developed below. In particular, when the discount factor is generated by the short rate, the Feynman--Kac formula yields the PDE representation of arbitrage-free bond prices.





% The old Neural Networks as Function Approximators:


\section{Neural Networks as Function Approximators}

A feedforward neural network defines a parametric family of functions. The goal is to find parameters $\Omega$ such that $f(x;\Omega)$ approximate some target function $f^*$ well 

%that minimise the discrepancy between $f(x;\Omega)$ and $f^*(x)$ over the training data.

%The architecture of a network, meaning the number of layers, the width of each layer, and the activation functions, is fixed before training. 

The parameters $\Omega = \{W_h, b_h\}_{h=1}^H$ are then free, and each choice of $\Omega$ produces a different function. The full set of functions the network can represent is therefore the class
\[
    \bigl\{f(\,\cdot\,;\Omega) : \Omega \in \Theta\bigr\},
\]
where $\Theta$ denotes the parameter space. Equivalently, after identifying all weights and biases with a vector of parameters, one may regard $\Theta \subseteq \mathbb{R}^p$, where $p$ is the total number of parameters.

Training is the process of searching through this class to find a parameter choice $\Omega$ for which $f(\,\cdot\,;\Omega)$ approximates $f^*$ well on the available data. 




A natural question is whether this class is rich enough to approximate any target function of interest, or whether the architecture imposes fundamental limitations. The following theorem gives the classical finite-dimensional universal approximation property. We state it in notation consistent with the neural-network framework introduced above. A classical result of this type is due to \citet{Hornik1989}.

\begin{theorem}[Finite-dimensional universal approximation theorem, cf.\ \citealt{Hornik1989}]
\label{thm:UAT}
Let $\psi : \mathbb{R} \to \mathbb{R}$ be continuous and non-constant. Then for any compact set $K \subset \mathbb{R}^{M_0}$, any continuous target function $f^* : K \to \mathbb{R}$, and any $\varepsilon > 0$, there exist a width $M_1 \in \mathbb{N}$ and parameters
\[
    \Omega = (W_1,b_1,W_2,b_2),
\]
with
\[
    W_1 \in \mathbb{R}^{M_1 \times M_0},
    \qquad
    b_1 \in \mathbb{R}^{M_1},
    \qquad
    W_2 \in \mathbb{R}^{1 \times M_1},
    \qquad
    b_2 \in \mathbb{R},
\]
such that the network
\[
    f(x;\Omega) = W_2\,\psi(W_1x+b_1) + b_2,
\]
where $\psi$ is applied element-wise, satisfies
\[
    \sup_{x \in K} |f^*(x)-f(x;\Omega)| < \varepsilon.
\]
\end{theorem}

The approximating class in Theorem~\ref{thm:UAT} is a feedforward network with a single hidden layer of width $M_1$ and an identity activation in the output layer, and is therefore a special case of the broader architecture in \eqref{eq:network}. The theorem shows that even this restricted architecture can approximate any continuous scalar-valued function on a compact set to arbitrary accuracy.

However, the theorem is an existence result only. It does not specify how large the width $M_1$ must be, whether deeper architectures can achieve the same approximation more efficiently, or how suitable parameters are to be found in practice. Thus, for a fixed architecture, approximation may still be poor if the network is too narrow, and suitable parameters may remain difficult to obtain through training. In practice, such questions of architecture design and optimisation are therefore addressed empirically, typically through model selection and validation on data not used for training.